\section{Setting and notation}\label{sec:background}

\subsection{Bias and Cayley graphs}
A \emph{weighted sample space} is a finite set $S\subset\N$ with weights $w:S\to[0,\infty)$ and mass
$M=\sum_{n\in S}w(n)>0$. For a Dirichlet character $\chi$ modulo $q$ put
\[
  \bias_\chi(S,w)=\frac{\bigl\lvert\sum_{n\in S}w(n)\chi(n)\bigr\rvert}{M}.
\]
If every $n\in S$ is coprime to $q$, then $(S,w)$ defines the probability measure $\mu(g)=M^{-1}\sum_{n\equiv
g}w(n)$ on $\Gq$, and $\bias_\chi(S,w)=\abs{\mathbb E_\mu\chi}$.

\begin{lemma}[eigenvalues; \lean{GRHPseudo.charSum_eq_eigenvalue}]\label{lem:eigen}
Let $T f(g)=\sum_{n\in S}w(n)f(gn)$ for $f:\Z/q\Z\to\C$. For every Dirichlet character $\chi$ modulo
$q$ one has $T\chi=\bigl(\sum_{n\in S}w(n)\chi(n)\bigr)\chi$.
\end{lemma}
\begin{proof}
$T\chi(g)=\sum_n w(n)\chi(g)\chi(n)$, since $\chi$ is completely multiplicative.
\end{proof}

Since the characters of $\Gq$ form an orthogonal basis of functions on $\Gq$, the bias of $\mu$ equals the
second-largest eigenvalue modulus of the normalized weighted Cayley graph $\mathrm{Cay}(\Gq,\mu)$. We
formalize only the eigenvector identity; the completeness statement is standard.

\subsection{Dirichlet $L$-functions and the GRH theorem}
For $\chi$ modulo $q$ we write $L(s,\chi)$ for Mathlib's \lean{DirichletCharacter.LFunction}. It is the
analytic continuation of $\sum_{n\ge1}\chi(n)n^{-s}$ (with $\chi(n)=0$ for $\gcd(n,q)>1$), and it is entire
for $\chi\ne\chi_0$. The GRH theorem we use is the compact edition's \lean{LFunction_ne_zero_of_half_lt_re}
\cite{grhlean}:
\begin{equation}\label{eq:grh}
  \forall q\ge1,\ \forall\chi \bmod q,\ \forall s,\quad \operatorname{Re}s>\tfrac12,\ (\chi\ne\chi_0\text{ or } s\ne1)
  \ \Longrightarrow\ L(s,\chi)\ne0 .
\end{equation}
We also use its consequences \lean{LFunction_eq_zero_re} (zeros with $0<\operatorname{Re}s<1$ lie on the line)
and \lean{riemannHypothesis} (Mathlib's \lean{RiemannHypothesis}).

\subsection{Fourier conventions and the explicit formula}
For $f\in L^1(\R)$ and $z\in\C$ put $\hat f(z)=\int_\R f(u)e^{izu}\,du$ whenever the integral converges, and
$\tilde f(u)=\overline{f(-u)}$. For $f\in C^2_c(\R)$ the Weil test function is $k=f*\tilde f$, i.e.\
$k(u)=\int f(t)\overline{f(t-u)}\,dt$. Then $\hat k(r)=\abs{\hat f(r)}^2\ge0$ for real $r$, and
$k(-u)=\overline{k(u)}$. For a zero $\rho$ we write $\gamma_\rho=(\rho-\tfrac12)/i$, which is real iff
$\operatorname{Re}\rho=\tfrac12$.

We use the explicit formula in the following form, proved in Lean in \cite{zeta23}. For $\chi$ primitive
modulo $q>1$ of parity $\kappa\in\{0,1\}$ and $k\in C^2_c(\R)$, with $m_\rho$ the multiplicity of the
nontrivial zero $\rho$, the zero sum converges absolutely and
\begin{multline}\label{eq:EFchi}
  \sum_\rho m_\rho\,\hat k(\gamma_\rho)
  = -\sum_{n\ge1}\frac{\Lambda(n)}{\sqrt n}\Bigl(\chi(n)k(\log n)+\overline{\chi(n)}k(-\log n)\Bigr)\\
  +\frac1{2\pi}\int_\R \hat k(r)\Bigl(\operatorname{Re}\psi\bigl(\tfrac14+\tfrac\kappa2+\tfrac{ir}2\bigr)+\log\tfrac q\pi\Bigr)dr .
\end{multline}
We write $-P_\chi(k)+\Gamma_{q,\kappa}(k)$ for the right side. For $\zeta$,
\begin{equation}\label{eq:EFzeta}
  \sum_\rho m_\rho\,\hat k(\gamma_\rho)=\hat k(\tfrac i2)+\hat k(-\tfrac i2)
  -\sum_{n\ge1}\frac{\Lambda(n)}{\sqrt n}\bigl(k(\log n)+k(-\log n)\bigr)
  +\frac1{2\pi}\int_\R\hat k(r)\bigl(\operatorname{Re}\psi(\tfrac14+\tfrac{ir}2)-\log\pi\bigr)dr ,
\end{equation}
which we write as $\mathrm{Pole}(k)-P_\zeta(k)+\Gamma_\zeta(k)$. The Lean statements are
\lean{Zeta23.ThmE.EF_lit_chi_L} and \lean{Zeta23.WeilEF.EF_lit_zetaZeroConfig}; their right sides are the
definitions \lean{literatureRHSchi} and \lean{literatureRHS}.

\subsection{The test function}
Let $S:\R\to[0,1]$ be the smooth transition function, with $S=0$ on $(-\infty,0]$, $S=1$ on $[1,\infty)$
and $S\in C^\infty$ (Mathlib's \lean{Real.smoothTransition}). Following \cite{grhmiller}, put
\[
  \eta_a(t)=S\bigl(64(a-t)\bigr)\,S\bigl(64(a+t)\bigr),\qquad f_a(t)=\cosh(t/2)\,\eta_a(t),\qquad k_a=f_a*\tilde f_a .
\]
Thus $0\le f_a\le\cosh(t/2)$, $f_a(t)=\cosh(t/2)$ for $\abs t\le a-\tfrac1{64}$, and $f_a=0$ for $\abs t\ge a$.
The kernel $k_a$ is real, nonnegative and even, and vanishes for $\abs u>2a$. We write $a'=a-\frac1{64}$ and
\[
  F_\pm(a)=\int_\R f_a(t)e^{\pm t/2}\,dt ,\qquad A(a)=\hat k_a(-\tfrac i2)=F_+(a)F_-(a),\qquad
  B(a)=\hat k_a(\tfrac i2)+\hat k_a(-\tfrac i2)=2F_+F_- .
\]
